% Section: Introduction

A motor-imagery (MI) brain--computer interface (BCI) must decide
\emph{what} the user intends and \emph{when} to commit that decision.
Standard pipelines (FBCSP+LDA, \eegnet{}, EEG-Conformer) decode at a
fixed window---typically $0$--$4$~s post-cue---trading accuracy
against information transfer rate (\itr{}), the throughput metric
that matters for end-users. Dynamic-stopping classifiers in the
style of the sequential probability ratio test
(SPRT)~\cite{Liu2017MotorImageryBTSPRT} introduce
evidence-based early termination, but these are hand-tuned
thresholds on confidence
statistics, not learned policies.

Recent
work~\cite{Aung2024EEG_RL_Net,Fidencio2025ErrPRL}
frames BCI decoding as sequential decision-making, but none use
offline reinforcement learning (RL) --- learning a policy purely
from logged data, with no online interaction --- none expose an
explicit \textsc{Defer} action as part of the learned policy, none
address motor imagery at scale, and---critically---none employ a
calibrated off-policy evaluator. Off-policy evaluation (OPE)
estimates the value a candidate policy \emph{would} obtain if
deployed, using only logged data; it is the tool that makes
offline policy selection auditable. Existing BCI evaluations
simulate a single classifier under a fixed protocol, which is
on-policy simulation in a deterministic environment, not OPE in
any rigorous sense.
Meanwhile, offline RL with calibrated OPE has matured in adjacent
clinical domains (sepsis, ventilator weaning, anesthesia depth
control)~\cite{Kumar2020CQL,Le2019BatchRL},
but this machinery has not been ported to BCI.

This paper fills that methodological gap. Our central
contribution is a single reframing with two halves: (a)~we turn
within-trial MI EEG decoding into an offline RL problem with a
clinically meaningful action space $\{\textsc{Commit}_k,
\textsc{Defer}, \textsc{Recal}, \textsc{Abstain}\}$, training a
Conservative Q-Learning (\cql{})~\cite{Kumar2020CQL} agent on
logged trajectories from two behavior policies; and (b)~we make
the resulting policies \emph{selectable offline} by validating a
fitted-Q evaluation (\fqe{}) estimator against on-policy
Monte-Carlo ground truth. Fig.~\ref{fig:system_overview} gives a
high-level view of the full system: the decision loop the agent
runs within each trial (top) and the offline training-and-selection
pipeline wrapped around it (bottom). Everything downstream ---
risk-sensitive critics, safety constraints, encoder swaps,
cross-subject benchmarks --- instantiates or stress-tests this
one reframing.

\paragraph*{Contributions}
(i)~\textbf{Calibrated OPE for BCI:} per-target-policy \fqe{}
attains Pearson $r\!=\!0.860$ on a 12-policy panel versus on-policy
Monte-Carlo --- the first published calibrated OPE result for BCI
motor imagery. (ii)~\textbf{Selective offline-RL decoder:} a
conditional-value-at-risk (CVaR) critic combined with a
constrained-MDP (CMDP) commit-error budget reaches commit
accuracy $87.6\!\pm\!2.9\%$ at $37.0$~bits/min~\itr{} on the
canonical 3-subject \bcia{} panel (EEG-Conformer encoder) and
$74.6\!\pm\!11.9\%$ at $24.9$~b/m on the full 9-subject panel
(Table~\ref{tab:sota}). These are \emph{selective} operating
points --- the agent commits on only $27\%$/$29\%$ of trials,
deferring or abstaining on the rest, so commit accuracy is not
comparable to the task accuracy of a classifier forced to answer
every trial (\S\ref{sec:results:sota} reports both) --- and this
is the first offline-RL BCI decoder with \textsc{Defer} as a
\emph{learned} action (distinct from post-hoc
thresholding~\cite{Ganeshkumar2017RejectMI}).
(iii)~\textbf{Cross-subject leave-one-subject-out (LOSO)
benchmarks on three datasets:} 62-channel
Lee2019 ($N\!=\!10$) yields commit accuracy
$0.696\!\pm\!0.104$ vs random $0.481$, $10/10$ wins, paired
Wilcoxon $p\!=\!0.001$ --- the first published Lee2019 LOSO
benchmark for offline-RL BCI, completing a monotonic
channel-count to LOSO-significance trend
($p\!=\!0.10$/$0.15$/$0.001$ on $\{3, 22, 62\}$-channel datasets).
Against an EEGNet-fixed mandatory-commit baseline on the same
splits, raw commit accuracy is statistically comparable
($p\!=\!0.43$); the agent's advantage is a $-25\%$ relative
reduction in wrong-commit rate ($0.214$ vs $0.284$,
$p\!=\!0.07$), positioning the contribution as a safety-oriented
selective decoder rather than an accuracy SOTA.
(iv)~\textbf{Foundation-model RL fine-tuning:} first end-to-end
offline-RL fine-tuning of the \labram{}~\cite{LaBraM2024} EEG
foundation model on a BCI task; ties \eegnet{} in aggregate
(paired Wilcoxon $p\!=\!0.79$) with a reproducible subject-specific
encoder-rescue effect ($+16.8$pp on \bcib{}, $+15.9$pp on \bcia{}
for the same held-out individual).
(v)~\textbf{Open release} of the entire pipeline (data, MDP
simulator, agent, OPE estimators, \labram{} integration) under the
MIT license.

\paragraph*{Topic match}
The work falls in two named topical areas of the venue:
\emph{Adaptive Control} and \emph{Neural Signal Processing}. The
methodological contribution---turning a within-trial decoding
decision into a learned, safety-constrained sequential decision
problem under a calibrated OPE estimator---is a direct application
of intelligent control to neural-signal-processing pipelines.
